\chapter{Foundations for Absolute Beginners}
\label{ch:foundations}

This chapter builds the two backgrounds you need before any hand tracking can make sense: the
Python language (Section~\ref{sec:python}) and how pictures live inside a computer
(Section~\ref{sec:images}). It closes with a plain-English account of what machine learning actually
is (Section~\ref{sec:ml}). Nothing here is specific to hands yet --- that is deliberate. If you have
never written a line of code, read every word. If you have written Python before, read the boxes
anyway: they explain the \emph{why} behind idioms you may already use by habit.

\section{Python from Scratch for Computer Vision}
\label{sec:python}

\subsection{What is an interpreter, and how does Python run a file?}

When you type \code{python hand_tracker.py} and press Enter, you are not ``starting a program'' in
the way you start a compiled \code{.exe}. You are starting a program called the \emph{interpreter},
and handing it a text file.

\begin{pythonconcept}{The interpreter, and why it matters here}
An \textbf{interpreter} is a program that reads your text file and performs the instructions one
statement at a time, top to bottom, translating as it goes.

A useful analogy: a \textbf{compiled} language is like translating an entire novel into another
language and then publishing it --- slow to prepare, fast to read. An \textbf{interpreted} language
is like a simultaneous translator standing next to you --- ready instantly, but working sentence by
sentence.

Two practical consequences you will meet within the first five minutes:

\begin{itemize}
  \item \textbf{Errors appear only when that line is reached.} If line 300 has a typo but your
        camera never opens, the program may die at line 40 and you will never see the line-300
        problem. An interpreter does not read the whole file first.
  \item \textbf{The same file can behave differently on different runs} if it depends on something
        outside the file (a camera, a network, a random number). This is why the tracker logs
        ``Camera 0 opened at 640x480'' --- so you know what the run actually did.
\end{itemize}

Python files are named with the \code{.py} suffix, and \code{python} is the interpreter. In this
project there are two of them: \code{hand_logic.py} and \code{hand_tracker.py}.
\end{pythonconcept}

\subsection{Variables and dynamic typing}

A variable is a name attached to a value. That is the whole idea.

\begin{codecard}
\begin{lstlisting}
fps = 28            # an int  (whole number)
alpha = 0.65        # a float (decimal number)
label = "Right"     # a str  (text)
is_up = True        # a bool (True / False)
\end{lstlisting}
\end{codecard}

These four lines use the four basic types you will meet constantly in this project. Python calls
its basic types \emph{built-in types}, and the project uses exactly these:

\begin{center}
\rowcolors{2}{slatebg}{white}
\begin{tabularx}{\textwidth}{@{}l l X@{}}
\toprule
\rowcolor{white}
\textbf{Type} & \textbf{Example} & \textbf{Used in the project for} \\
\midrule
\code{int}   & \code{320}      & pixel coordinates, frame counters, finger counts \\
\code{float} & \code{0.65}     & normalized coordinates, the palm normal $Z$, alpha values \\
\code{str}   & \code{"Right"}  & handedness labels, gesture names, file paths \\
\code{bool}  & \code{True}     & ``is this finger extended?'' --- five per hand, every frame \\
\bottomrule
\end{tabularx}
\end{center}

\begin{pythonconcept}{Dynamic typing: the name has a type, not the box}
In C or Java you would write \code{int fps = 28;} --- the \emph{container} is declared as an integer
and can never hold anything else.

Python is \textbf{dynamically typed}: the \emph{name} simply points at a value, and the value knows
its own type. This line is legal:

\begin{lstlisting}
value = 5          # value points at an int
value = "hello"    # now the same name points at a string
\end{lstlisting}

The price of this freedom is that type mistakes are discovered at \emph{run time}, not
compile time. The project pays that price back with \textbf{type hints} and \textbf{tests} --- see
Section~\ref{sec:typehints}.
\end{pythonconcept}

\subsection{Lists, tuples and dictionaries}

Nearly every bug in a hand tracker is really a data-shape bug: the wrong number of points, points in
the wrong order, or a name that does not exist. So it is worth being precise about the three
containers this project uses.

\subsubsection*{Lists --- ordered, changeable}

A \textbf{list} is an ordered sequence you can modify. Square brackets, comma separated.

\begin{codecard}
\begin{lstlisting}
FINGER_PAIRS = [(8, 6), (12, 10), (16, 14), (20, 18)]
states = [False] * 5          # five False values
states[0] = True              # now [True, False, False, False, False]
print(len(states))            # 5
\end{lstlisting}
\end{codecard}

Line 1 is taken from \code{hand_logic.py}: it is the list of (fingertip, knuckle) landmark index
pairs for the four long fingers. Line 2 is the ``safe result'' the project falls back to whenever
landmarks are missing --- note that multiplying a list by an integer \emph{repeats} it, it does not do
arithmetic.

The most important property of a list here is \textbf{indexing}: \code{points[0]} is the first
element, \code{points[8]} is the ninth. In this project, landmark indices are the entire language we
speak --- index 0 is the wrist, index 8 is the tip of the index finger, index 17 is the base of the
pinky. If you remember one thing from this section, remember that \emph{the index number is a name}.

\subsubsection*{Tuples --- ordered, frozen}

A \textbf{tuple} is like a list that cannot be changed. Round brackets.

\begin{codecard}
\begin{lstlisting}
tip = (320, 240)          # a single point: x, y
x, y = tip                # unpacking: x becomes 320, y becomes 240
pixel = points[8]         # a tuple from the landmark list
dist = math.hypot(pixel[0] - other[0], pixel[1] - other[1])
\end{lstlisting}
\end{codecard}

A 2D point is naturally a tuple: it is a fixed pair, and it should not be accidentally mutated. The
project uses tuples for coordinates and lists for collections of coordinates, so the shape of the
data tells you what it is.

\subsubsection*{Dictionaries --- named, lookup by key}

A \textbf{dictionary} maps keys to values. Curly braces.

\begin{codecard}
\begin{lstlisting}
hand = {
    "label": "Right",
    "px": [(320, 420), (365, 400)],   # pixel landmarks
    "palm_z": -8775.0,
    "orientation": "Palm",
}

print(hand["label"])        # Right
print(hand["bbox"])         # KeyError!  this key does not exist
\end{lstlisting}
\end{codecard}

This is precisely the structure \code{parse_hand_data} builds once per frame for each detected hand.
A dictionary is the right choice because a hand has \emph{named} properties whose order does not
matter --- and because a reader of the code can ask for \code{hand["orientation"]} instead of
remembering ``the fifth element''.

\begin{edgetrap}{AttributeError and KeyError --- the two errors of this section}
\begin{itemize}
  \item \code{hand["bboxx"]} raises \code{KeyError}: the name is not in the dictionary.
  \item \code{hand.label} raises \code{AttributeError}: dictionaries are indexed with
        \code{[]}, not with a dot. The dot syntax is for \emph{objects} (Section~\ref{sec:oop}).
\end{itemize}
These two messages account for a large share of real-world Python errors in computer-vision code,
because MediaPipe results are objects while our cached hand data is a dictionary --- the access
syntax differs between them.
\end{edgetrap}

\subsection{Functions, arguments, return values and type hints}
\label{sec:typehints}

A \textbf{function} is a named block of work with inputs and (usually) one output.

\begin{codecard}
\begin{lstlisting}
def euclidean_distance(a, b):
    return math.hypot(a[0] - b[0], a[1] - b[1])

gap = euclidean_distance((0, 0), (3, 4))    # gap is 5.0
\end{lstlisting}
\end{codecard}

Reading that from left to right:

\begin{itemize}
  \item \code{def} announces a function definition.
  \item \code{euclidean\_distance} is its name.
  \item \code{(a, b)} are its \textbf{parameters} --- placeholders for the values it will receive.
  \item The \textbf{indented block} is the \textbf{body}.
  \item \code{return} hands a value back to whoever called it. A function without a \code{return}
        hands back a special value called \code{None}, meaning ``nothing''.
\end{itemize}

\subsubsection*{Default arguments}

\begin{codecard}
\begin{lstlisting}
def fingers_state(pixels, hand_label="Unknown", thumb_mode="distance"):
    ...

fingers_state(points)                       # label is "Unknown", mode "distance"
fingers_state(points, "Right")              # label is "Right"
fingers_state(points, thumb_mode="x")       # keyword argument: order does not matter
\end{lstlisting}
\end{codecard}

An argument with an \code{=} in the definition is \textbf{optional}. This is why the project's
public methods can be called simply --- \code{tracker.find\_positions(frame, 0)} --- while still
offering knobs such as \code{alpha=0.65}.

\subsubsection*{Type hints are documentation that can be checked}

\begin{codecard}
\begin{lstlisting}
def to_pixel_landmarks(landmarks, width, height):
    """Convert normalized landmarks to integer pixel coordinates."""
    return [(int(lm.x * width), int(lm.y * height)) for lm in landmarks]
\end{lstlisting}
\end{codecard}

Python does not enforce the types written in a hint --- it will happily run and crash later if you
pass nonsense. The hint is a promise to the reader, and modern editors check those promises for you
as you type, before you ever run the program. Throughout this handbook you will see hints such as
\code{list[bool]} (a list of booleans), \code{tuple[int, int]} (a pair of integers --- one pixel
coordinate) and \code{list[dict]} (one dictionary per detected hand).

\begin{keyidea}
\textbf{Reading a signature aloud} is a genuinely useful skill. Take
\code{fingers\_state(pixel\_landmarks, hand\_label="Unknown", palm\_z=None, thumb\_mode="distance")}.
Read as: ``give me the pixel landmarks; optionally tell me which hand this is; give me five booleans
back --- one per finger, thumb first.'' Once you can read signatures, the architecture of the whole
project becomes legible.
\end{keyidea}

\subsection{Object-oriented programming, simplified}
\label{sec:oop}

Everything so far has been \emph{data} (values) and \emph{procedures} (functions) kept apart. Object
orientation glues them together.

\begin{pythonconcept}{Class, object, self --- the three words that block beginners}
A \textbf{class} is a blueprint. An \textbf{object} (or \emph{instance}) is a thing built from that
blueprint. \code{self} is how the object refers to itself.

Analogy: a class is the architectural drawing of a house; an object is an actual house. You can build
many houses from one drawing, and each has its own front door colour --- its own data --- while
sharing the same structure.

\begin{lstlisting}
class HandTracker:
    def __init__(self, max_hands=2):
        self.max_hands = max_hands     # store on THIS object
        self.hands_data = []           # each object gets its own empty list

    def hand_count(self):
        return len(self.hands_data)

tracker = HandTracker()                # build an object from the blueprint
print(tracker.hand_count())            # 0
\end{lstlisting}

Line by line:

\begin{itemize}
  \item \code{class HandTracker:} defines the blueprint. By convention, class names are capitalised.
  \item \code{\_\_init\_\_} is the \textbf{constructor}: it runs automatically when the object is
        built. The name has two underscores on each side --- in Python that signals ``special method
        that the language itself calls''.
  \item \code{self} is the first parameter of every method. Python passes the object in
        automatically, so \code{tracker.hand\_count()} is really
        \code{HandTracker.hand\_count(tracker)}. You never pass \code{self} yourself; you only
        declare it.
  \item \code{self.hands\_data = []} creates a piece of state \emph{belonging to this object}.
        Crucially, each object gets its own list. Two trackers would not share one hand cache.
\end{itemize}

Why does this project need a class at all? Because detection has \textbf{state}: the MediaPipe model
object itself, the smoothing history, the gesture debouncers. A bare function would have to be handed
all of that on every call. A class lets the state live next to the methods that use it --- which is
why the detector is \code{HandTracker}, and \code{find\_hands}, \code{find\_positions} and
\code{fingers\_up} are its methods.
\end{pythonconcept}

\subsection{Five special constructs used in this project}

\subsubsection*{1. \code{deque} with a maximum length}

\begin{codecard}
\begin{lstlisting}
from collections import deque
from hand_logic import GestureDebouncer

history = deque(maxlen=5)
history.append("Fist")        # ['Fist']
history.append("Peace")       # ['Fist', 'Peace']
# ... after five appends ...
history.append("Fist")        # the OLDEST entry is silently evicted
\end{lstlisting}
\end{codecard}

A \code{deque} (``deck'', double-ended queue) is a list optimised for adding and removing at both
ends. The \code{maxlen} argument makes it a \textbf{sliding window}: it never grows beyond five
entries, and it never needs to be trimmed by hand. That single feature is the entire implementation of
the gesture debouncer in \code{hand_logic.py}: keep the last five gestures, count how many agree.

\begin{pythonconcept}{Why not a list?}
You could write \code{history.append(g); if len(history) > 5: history.pop(0)} --- but
\code{pop(0)} on a list shifts every remaining element, which is $O(n)$. A \code{deque} is designed
for that operation and evicts in $O(1)$. It is also harder to forget. In a 30-frames-per-second
loop, forgetting to trim a list is a slow memory leak that only shows up during a long demo.
\end{pythonconcept}

\subsubsection*{2. List comprehensions}

\begin{codecard}
\begin{lstlisting}
# the long way
pixels = []
for lm in landmarks:
    pixels.append((int(lm.x * width), int(lm.y * height)))

# the comprehension
pixels = [(int(lm.x * width), int(lm.y * height)) for lm in landmarks]
\end{lstlisting}
\end{codecard}

Both blocks do the same thing. The second is the idiom you will see in this codebase wherever one
collection is transformed into another. Read a comprehension as: ``build this expression, for each
item in this collection.''

\subsubsection*{3. \code{if \_\_name\_\_ == "\_\_main\_\_":}}

\begin{codecard}
\begin{lstlisting}
def main():
    ...

if __name__ == "__main__":
    raise SystemExit(main())
\end{lstlisting}
\end{codecard}

\begin{pythonconcept}{The most-asked Python question, answered properly}
When Python runs a file \emph{directly} (\code{python hand\_tracker.py}), it sets a built-in
variable called \code{\_\_name\_\_} to the string \code{"\_\_main\_\_"} inside that file.

When the same file is \emph{imported} by another file (\code{import hand\_tracker}), Python sets
\code{\_\_name\_\_} to the module's own name, \code{"hand\_tracker"}.

So \code{if \_\_name\_\_ == "\_\_main\_\_":} means \textbf{``run this block only if I am the program
being started, not if I am being used as a library.''}

This project needs exactly that distinction for three reasons:

\begin{itemize}
  \item The test suite imports \code{hand\_tracker} to test its helper functions. If the capture
        loop ran at import time, \code{pytest} would try to open your webcam.
  \item \code{raise SystemExit(main())} converts the integer that \code{main()} returns into the
        process exit code, so CI can tell success (0) from failure (1).
  \item It keeps the top of the file importable for tooling without side effects.
\end{itemize}
\end{pythonconcept}

\subsubsection*{4. Exceptions and \code{finally}}

\begin{codecard}
\begin{lstlisting}
try:
    while True:
        ok, frame = capture.read()
        ...
finally:
    capture.release()          # ALWAYS runs: normal exit, error, or Ctrl+C
    cv2.destroyAllWindows()
\end{lstlisting}
\end{codecard}

An \textbf{exception} is Python's way of saying ``I cannot continue''. \code{try} marks a region
where you expect one; \code{finally} marks cleanup that must happen no matter how the region ends.
The tracker's capture loop uses \code{finally} to release the camera and close the window. Without
it, an unplugged webcam or a \code{Ctrl+C} would leave the camera device locked, and the next run
would fail with ``could not open a webcam'' until you rebooted --- a classic and infuriating bug.

\subsubsection*{5. The walrus-free guard clause}

\begin{codecard}
\begin{lstlisting}
def find_positions(self, frame, hand_index=0):
    if 0 <= hand_index < len(self.hands_data):
        return list(self.hands_data[hand_index]["px"])
    return []
\end{lstlisting}
\end{codecard}

Note the \emph{chained comparison} \code{0 <= i < len(...)}. Python evaluates it as
\code{(0 <= i) and (i < len(...))}, which is why a negative index --- which would otherwise mean
``count from the end'' --- is rejected. This two-line pattern (``check, return early; otherwise fall
through to the safe default'') is how the project removes an entire class of crashes:
``IndexError: list index out of range'' is impossible here.

\section{How Digital Images Work}
\label{sec:images}

\subsection{Pixels, resolution and channels}

An image is a grid of coloured dots called \textbf{pixels} (short for ``picture elements''). The
grid has a width and a height, measured in pixels; together they are the \textbf{resolution}.

\begin{mlconcept}{From photons to numbers}
A camera sensor measures, for each of its tiny light-sensitive cells, how much light arrived. The
webcam in this project is asked for $640 \times 480$ cells (that is 307{,}200 pixels). For each cell
we get three intensities: how much \textbf{red}, how much \textbf{green}, how much \textbf{blue} light
arrived.

Each intensity is stored as one byte --- an integer from 0 (none) to 255 (full). Three bytes per
pixel, 307{,}200 pixels: about 0.9 megabytes per frame, and at 30 frames per second roughly
27 MB/s flowing through your program. This is exactly why the project caches work per frame
instead of recomputing it (Chapter~\ref{ch:architecture}).

A colour is therefore an ordered triple. \code{(255, 0, 0)} is pure red in RGB order;
\code{(0, 0, 0)} is black; \code{(255, 255, 255)} is white.
\end{mlconcept}

\subsection{Why OpenCV thinks in BGR}

Here is a fact that costs beginners hours: OpenCV --- the library that reads your webcam --- stores
colour channels in the order \textbf{blue, green, red}, not red, green, blue. It is a historical
inheritance from the 1990s BGR camera hardware OpenCV was written against.

MediaPipe, on the other hand, was trained on ordinary RGB images.

\begin{edgetrap}{The BGR/RGB mismatch: silently wrong, not loudly broken}
If you feed a BGR image to an RGB model, nothing crashes. Every colour is simply permuted: red and
blue swap. Skin becomes a strange blue-ish tint, and the detector's accuracy quietly degrades --- or
drops to zero --- on some skin tones and lighting conditions.

It is the worst kind of bug: no exception, no error message, just worse results. This is why the
pipeline contains exactly one explicit conversion, at exactly one place:

\begin{lstlisting}
rgb = cv2.cvtColor(frame, cv2.COLOR_BGR2RGB)
results = self.hands.process(rgb)
\end{lstlisting}

Everything {\em before} that line is BGR (because it came from the camera and will be shown by
\code{cv2.imshow}); everything {\em after} is MediaPipe's business. Screen output stays in BGR,
because \code{cv2.imshow} expects BGR. The conversion is a one-way door for the model only.
\end{edgetrap}

\subsection{A frame is a three-dimensional array}

This is the single most useful mental model in the whole project.

\begin{mlconcept}{The flipbook and the cube}
Think of video as a \textbf{flipbook}: each page is one still image, and flipping fast enough
produces motion. A single frame is one page.

Now zoom into one page. It is a grid of pixels. Add the three colour values per pixel and the page
becomes a \textbf{three-dimensional block of numbers}:

\begin{itemize}
  \item \textbf{axis 0} --- the rows (height): $480$ of them
  \item \textbf{axis 1} --- the columns (width): $640$ of them
  \item \textbf{axis 2} --- the channels: $3$ of them
\end{itemize}

In NumPy --- Python's numerical array library --- a frame is described by its \code{shape}, and you
can inspect it at any moment:

\begin{lstlisting}
>>> frame.shape
(480, 640, 3)
>>> frame[0, 0]          # top-left pixel, all three channels
array([ 42,  42,  45], dtype=uint8)
>>> frame[240, 320]      # the pixel at the centre of the frame
array([118, 147, 201], dtype=uint8)
>>> frame[240, 320, 0]   # just the BLUE value of that pixel
118
\end{lstlisting}

Note the order: \code{frame[y, x]}, \emph{row first, column second}. This is a genuine trap, and it
is the reason the project writes \code{height, width = frame.shape[:2]} rather than the other way
around. MediaPipe works in the opposite convention --- it reports normalized $(x, y)$ --- so every
conversion in this codebase is a place where the two conventions meet.
\end{mlconcept}

\begin{keyidea}
\textbf{One sentence to memorise:} a frame is a \code{(height, width, 3)} block of bytes, indexed
\code{[row, column, channel]}, and a landmark is a normalised \code{(x, y)} pair indexed the other
way round. Most ``my dot is in the wrong place'' bugs are one of these two facts being forgotten.
\end{keyidea}

\section{Machine Learning Demystified}
\label{sec:ml}

\subsection{Rules plus data, versus data plus answers}

\begin{mlconcept}{The inversion that defines machine learning}
\textbf{Traditional programming.} You write the rules, you feed in data, you get answers out.

\begin{center}
\code{rules + data $\rightarrow$ answers}
\end{center}

Example: ``if the pixel's blue value is greater than its red value, call it skin.'' You wrote the
rule. The computer applies it.

\textbf{Machine learning.} You supply data and the correct answers for that data. The computer
\emph{derives} the rules.

\begin{center}
\code{data + answers $\rightarrow$ rules}
\end{center}

Example: you show the computer 30{,}000 photographs of hands together with the exact pixel location
of every knuckle, and it works out its own rules for finding knuckles.

You cannot write the rule for ``where is the pinky knuckle'' by hand. There is no threshold on red,
green or blue that finds it. That is precisely why machine learning is used for this half of the
project --- and why the other half (deciding whether a finger is up) is ordinary geometry that
\emph{can} be written by hand, and therefore \emph{should} be.
\end{mlconcept}

\subsection{What a neural network actually is}

\begin{mlconcept}{Weights, layers, and feature detectors}
Strip away the vocabulary and a neural network is a very long chain of multiply-and-add operations,
with a non-linear squashing function between the steps. That is all.

\begin{itemize}
  \item A \textbf{weight} is one number that gets multiplied by an input. A small network has
        thousands; a large one has billions.
  \item A \textbf{layer} is a group of weights applied together, producing a new set of numbers from
        the previous set.
  \item A \textbf{feature} is a pattern that a layer reacts to. Early layers react to edges and
        colour gradients. Middle layers react to combinations such as ``a curved edge above a
        straight one''. Late layers react to whole structures: ``this cluster of edges is a
        knuckle''.
\end{itemize}

An analogy that holds up better than most: think of a factory production line. Raw pixels go in at
one end. At each station, workers check for one specific simple thing and pass on a summary. By the
end of the line, the summary is abstract enough to answer the question --- ``the pinky knuckle is at
$(0.31, 0.62)$''.

The important consequence for you: \textbf{each layer sees only the previous layer's summary}, never
the raw image. That is what makes the whole thing fast enough to run on a CPU at 30 FPS.
\end{mlconcept}

\subsection{Training versus inference}

\begin{mlconcept}{We only ever run the trained model}
\textbf{Training} is the process of searching for the weights. It requires a large labelled dataset,
a great deal of computation --- usually GPUs for hours or days --- and the ability to measure error
and correct it. Training is what Google did once, in order to ship MediaPipe.

\textbf{Inference} is using a network whose weights are already fixed. Given an input, produce an
output. No learning happens. It is a one-directional calculation.

This project is \textbf{100\% inference}. The weights arrive inside the MediaPipe wheel as a
\code{.tflite} file --- a few megabytes of numbers --- and the library loads them into memory when
you write:

\begin{lstlisting}
self.hands = mp_hands.Hands(
    static_image_mode=mode,
    max_num_hands=max_hands,
    min_detection_confidence=det_conf,
    min_tracking_confidence=track_conf,
)
\end{lstlisting}

Those four keyword arguments are \emph{not} training settings. They choose how confident the model
must be before it reports something. Chapter~\ref{ch:engine} covers what each one does and why the
defaults are $0.7$ and $0.6$.
\end{mlconcept}

\subsection{Normalized coordinates, and the conversion to pixels}

This is the last prerequisite, and it is where the first real mathematics of the project appears.

MediaPipe does not report pixel positions. It reports \textbf{normalized coordinates}: numbers in the
range $[0.0, 1.0]$, where $0.0$ is the left (or top) edge of the image and $1.0$ is the right (or
bottom) edge.

\begin{mathdeepdive}{From normalized coordinates to pixels}
\textbf{Why normalized?} Because the model must not care about resolution. Whether the frame is
$640 \times 480$ or $1920 \times 1080$, a hand in the middle of the picture is at
$x_{\text{norm}} = 0.5$. Normalizing makes one trained model valid on every camera, forever. The
model predicts \emph{relative} position; the exact pixel is an application concern.

\textbf{The conversion.} If a landmark sits at normalized position $x_{\text{norm}}$ on an image
$W$ pixels wide, its pixel column is

$$x_{\text{px}} = x_{\text{norm}} \cdot W$$

and, since the image is $H$ pixels tall,

$$y_{\text{px}} = y_{\text{norm}} \cdot H$$

\textbf{Worked example.} MediaPipe reports the index fingertip at
$(x_{\text{norm}}, y_{\text{norm}}) = (0.5469,\ 0.4896)$ on the default $640 \times 480$ frame:

\begin{align*}
x_{\text{px}} &= 0.5469 \times 640 = 350.0 \\
y_{\text{px}} &= 0.4896 \times 480 = 235.0
\end{align*}

so the fingertip is drawn at pixel $(350, 235)$ --- measured from the \emph{top-left} corner.

\textbf{Why the code truncates instead of rounding.} The project writes

\begin{lstlisting}
return [(int(lm.x * width), int(lm.y * height)) for lm in landmarks]
\end{lstlisting}

and \code{int()} \emph{truncates} ($350.9 \rightarrow 350$) rather than rounds
($350.9 \rightarrow 351$). For drawing a filled circle of radius 10 pixels, a half-pixel bias is
invisible, and truncation is measurably faster. What matters is that the value is an integer at all:
OpenCV's drawing functions accept floats, but integer pixel coordinates make the cache smaller and
the arithmetic exact when the geometry chapter computes distances and cross products.
\end{mathdeepdive}

\begin{keyidea}
\textbf{Chapter summary.} Python is an interpreted language whose lists, tuples and dictionaries give
the project its data shapes. An image is a \code{(height, width, 3)} block of bytes in BGR order, and
MediaPipe speaks normalized RGB. Machine learning inverts the usual programming equation --- data
plus answers produces rules --- and this project uses a model that has already been trained, only
ever running inference. Landmarks arrive normalized and must be multiplied by the frame size to
become pixels.
\end{keyidea}
